\section{Research Design}

\textbf{Research question addressed:} How can geographic variation in LLM-generated mental health guidance be measured reliably across a diverse set of cities and income settings?

This dissertation uses a controlled audit design. A fixed battery of mental-health scenarios is presented to multiple large language models, with the stated city and country varied across queries while language and scenario wording are held constant. The design therefore estimates associations with the location stated in the prompt within the selected sample. It does not isolate which city-level attribute causes those differences, because income, WHO region, culture, language environment, healthcare infrastructure and crisis-service availability are correlated. Automated coding is applied to all 1,120 responses and compared with human labels on stratified samples. The revised actionability and localisation rules were first developed on 112 labels, frozen, and then evaluated on a fresh blinded 160-response hold-out. The hold-out gate required weighted $\kappa\geq0.70$ for ordinal outcomes and sensitivity, specificity and $F_1\geq0.80$ for binary outcomes.

\section{Dataset Construction}

\subsection{Factorial Structure}

The dataset is fully crossed: $7$ models $\times$ $20$ cities $\times$ $8$ scenarios $= 1{,}120$ responses, with exactly one response per cell. This balanced design ensures that differences between income groups or WHO regions are not confounded by unequal representation of models or scenarios, and that the effective sample size for any cross-city contrast is equal across all geographic groupings.

\subsection{City and Country Selection}

Twenty cities were selected to span four World Bank income categories (Low, Lower-Middle, Upper-Middle, High), with five cities per tier, and six WHO regions (AFR, AMR, EMR, EUR, SEAR, WPR). Cities were chosen to provide geographic and cultural diversity within each income tier rather than clustering in any single region. The full city list, income categories, WHO regions, and documented crisis-service categories are provided in Appendix B.

\subsection{Scenario Battery}

Eight mental health scenarios were constructed, spanning a range of acuity from moderate emotional distress (persistent loneliness, caregiver strain, sleep difficulties) through high-acuity presentations consistent with an implied risk of self-harm. Scenarios are reproduced in full in Appendix A. Each scenario was combined with a location statement of the form \textit{``I live in [city], [country]''} to create the prompt, with all other text identical across all twenty locations.

\subsection{Model Selection}

Seven models were queried: GPT-OSS 120B, Qwen 3 32B, Google Gemma 4 26B A4B Instruct, Meta Llama 4 Scout 17B-16E Instruct, Meta Llama 3.3 70B Instruct, Nemotron-3-Super-120B-A12B, and Mistral Small 2506. Responses were collected via each provider's API at a fixed temperature, with one sample per cell. A provenance audit confirmed that all 1,120 cells were filled with no missing, duplicated, or failed responses.

\subsection{Data Integrity}

A pre-analysis integrity audit identified that Nemotron-3-Super-120B-A12B returned truncated output in 143 of its 160 responses, consistent with a token ceiling applied during collection. Truncation was unrelated to city income category ($\chi^2$ test, $p = 0.99$) and is handled by including model identity as a fixed effect in all provisional hypothesis models. The response identifier supplied in the source files was unusable (missing in two files; populated with a different model's values in a third); it was reconstructed from the combination of city, scenario, and model name, verified to yield 1,120 unique keys.

\section{NLP Pipeline}

\textbf{What question this addresses:} How can the actionability, localisation, and cultural framing of free-text LLM responses be operationalised as numeric outcomes suitable for statistical modelling?

The primary NLP pipeline converts each free-text response into a numeric feature vector through a sequence of rule-based steps. No machine-learning model is used in this primary extraction pipeline: classifications are produced by fixed regular expressions, keyword vocabularies and explicit decision rules, making every decision traceable to a specific text pattern and auditable by inspection. The pipeline is implemented in Python using the \texttt{re} module and runs identically on all 1,120 responses. A separate prompt-based transformer experiment was evaluated on a 40-response subset as an alternative automated coder; it was not used to generate the outcomes for RQ1--RQ3.

\textbf{Why rule-based rather than embedding-based:} The rule-based approach was chosen because it makes every classification decision transparent and independently replicable from the response text alone. The known limitation is that it cannot capture paraphrase: a suggestion to ``speak with someone trained in mental wellness'' would not match the pattern \texttt{{\textbackslash}btherapist{\textbackslash}b}. This is a pre-specified limitation, documented in the Metric Specification Table (electronic supplementary material).

\subsection{Alternative Modern NLP Automated Coder}

As a secondary measurement-development experiment, the instruction-tuned transformer \texttt{Qwen/Qwen2.5-1.5B-Instruct} was tested as a few-shot, prompt-based content coder. This model is distinct from Qwen 3 32B, one of the seven models that generated the substantive study responses. Each response was combined with explicit coding definitions and two labelled examples. Qwen tokenised this input, represented the tokens through its learned contextual embeddings, processed them through transformer self-attention layers, and generated a JSON object containing seven labels: overall actionability, coping step, professional help, social support, crisis action, follow-up and surface localisation. No embedding vectors were extracted, no downstream classifier was trained, and no task-specific fine-tuning or parameter updating was performed. The method is therefore generative in-context classification, not an embedding-classifier pipeline.

The experiment was executed locally in Google Colab with deterministic decoding (\texttt{do\_sample=False}; recorded temperature 0.0). Generated records were checked for parseable JSON, the complete set of required fields, permitted integer values and logical dependencies between overall actionability and its component codes. The untouched Qwen predictions were compared with the frozen rule-based predictions only for the 40 common sample identifiers. Round-1 human labels served as an operational comparison reference because they were produced under the original coding framework; they were not treated as error-free ground truth. Direct comparisons were restricted to actionability, coping steps, professional help and surface localisation because the available rule-based outputs did not contain directly equivalent measures for social support, crisis action or follow-up.

\begin{figure}[H]
\centering
\begin{tikzpicture}[
  node distance=0.55cm,
  box/.style={draw, rounded corners, align=center, minimum width=3.4cm, minimum height=0.72cm, fill=blue!6},
  decision/.style={draw, rounded corners, align=center, minimum width=3.4cm, minimum height=0.72cm, fill=orange!10},
  arrow/.style={-{Latex[length=2mm]}, thick}
]
\node[box] (text) {Response text + frozen prompt};
\node[box, below=of text] (tokens) {Qwen tokenisation and contextual\\representations};
\node[box, below=of tokens] (generate) {Self-attention and JSON label\\generation};
\node[decision, below=of generate] (checks) {Parsing, value and dependency checks};
\node[box, below=of checks] (compare) {Matched comparison with rule-based\\and human-coded labels ($n=40$)};
\draw[arrow] (text) -- (tokens);
\draw[arrow] (tokens) -- (generate);
\draw[arrow] (generate) -- (checks);
\draw[arrow] (checks) -- (compare);
\end{tikzpicture}
\caption{Prompt-based Modern NLP automated-coding and validation workflow. Contextual embeddings are internal to Qwen; no vectors were extracted for a separate classifier.}
\label{fig:qwen_pipeline}
\end{figure}

This experiment was governed by the same validation principle as the primary pipeline: technical production of labels was insufficient for substantive use. Because the Qwen run showed weak agreement and internal dependency violations, it was rejected as a measurement instrument and retained only in the methodological development and reproducibility record.

\subsection{Pre-processing: Reasoning-Block Removal}

Qwen 3 32B wraps its chain-of-thought reasoning in \texttt{<think>...</think>} tags prior to its substantive response. Of the 1,120 responses, 160 --- all from Qwen 3 32B --- contained such blocks. If not removed, the internal reasoning would be treated as part of the visible response, inflating keyword counts. All \texttt{<think>} blocks are stripped before any other processing:

\begin{lstlisting}[language=Python]
THINK_BLOCK_RE = re.compile(
    r'<think>.*?</think>',
    flags=re.IGNORECASE | re.DOTALL
)
analysis_text = THINK_BLOCK_RE.sub(' ', response_text).strip()
\end{lstlisting}

The original \texttt{response\_text} field is preserved unchanged. All 160 blocks were properly terminated; the minimum post-stripping length was 300 words, confirming no response was rendered empty.

\subsection{Contact and Specificity Extraction}

Phone numbers are detected using a three-branch regular expression:

\begin{lstlisting}[language=Python]
PHONE_RE = re.compile(
    r'(?<!\d)(?:'
    r'\+\d{1,3}[\s\-]?(?:\(?\d{2,4}\)?[\s\-]?){1,4}\d{2,4}'
    r'|'
    r'(?:\(?\d{2,4}\)?[\s\-]){1,4}\d{2,4}'
    r'|'
    r'\d{6,}'
    r')(?!\d)')
\end{lstlisting}

Branch 1 matches international numbers beginning with a country code (\texttt{+}). Branch 2 matches domestic-format numbers with explicit separators. Branch 3 matches unspaced digit runs of six or more digits. A bare four-digit number (a year) matches none of the three branches. Recognised crisis short-codes (988, 116 123, 13 11 14, and others) are matched by a separate pattern \texttt{SHORT\_CODE\_RE} and their digit spans are masked before \texttt{PHONE\_RE} runs, preventing double-counting.

The specificity score $s_i$ for response $i$ is:

\begin{equation}
s_i = n^{\text{phone}}_i + n^{\text{short}}_i + n^{\text{url}}_i
\label{eq:specificity}
\end{equation}

where $n^{\text{phone}}_i$, $n^{\text{short}}_i$, and $n^{\text{url}}_i$ are the number of phone, short-code, and URL matches respectively. The binary indicator $\texttt{crisis\_number\_present}_i = \mathbf{1}[n^{\text{phone}}_i + n^{\text{short}}_i > 0]$.

\subsection{Keyword Vocabularies}

Five keyword vocabularies count mentions of support-type categories. Each vocabulary is a list of regular expressions applied case-insensitively:

\begin{equation}
\texttt{cat\_mentions}_i = \sum_{p \in V_{\text{cat}}} \lvert\{m : \text{re.findall}(p, t_i)\}\rvert
\label{eq:keyword_count}
\end{equation}

where $t_i$ is the \texttt{analysis\_text} for response $i$ and $V_{\text{cat}}$ is the vocabulary for category $\texttt{cat} \in \{\text{professional}, \text{family}, \text{community}, \text{religious}, \texttt{self\_management}\}$. The professional vocabulary contains nine patterns; the religious vocabulary contains 14 patterns including \texttt{{\textbackslash}bpray(?:er|ing|s)?{\textbackslash}b}, \texttt{{\textbackslash}bimam{\textbackslash}b}, and \texttt{{\textbackslash}bspiritual(?:ity)?{\textbackslash}b}.

\subsection{Directive-Context Counting}

Raw keyword counts do not distinguish recommendation from mention: ``many people avoid therapists'' and ``consider speaking with a therapist'' both increment the professional vocabulary count. The final pipeline counts keywords only when they co-occur within the same sentence or list item as a directive linguistic marker. The 29 directive markers include patterns such as \texttt{{\textbackslash}bconsider{\textbackslash}b}, \texttt{{\textbackslash}battend(?:ing)?{\textbackslash}b}, \texttt{{\textbackslash}byou (?:can|could|should){\textbackslash}b}, and \texttt{{\textbackslash}bI (?:would )?(?:recommend|suggest){\textbackslash}b}. An additional structural detector identifies bolded list-item headers as inherently directive.

\begin{equation}
\texttt{rec\_cat}_i = \sum_{s \in \text{sentences}(t_i)} \mathbf{1}[\exists\, d \in D : d \in s] \cdot \texttt{count}(V_{\text{cat}}, s)
\label{eq:directive_count}
\end{equation}

where $D$ is the set of directive markers and $\text{sentences}(t_i)$ is the sentence-level segmentation of $t_i$. \textbf{Decision:} this formulation replaced the original actionability outcome following validation failure ($F_1 = 0.680$). The rebuilt five-component measure achieved $F_1 = 0.744$ by restricting professional-referral detection to directive contexts, adding named coping steps, and reducing false positives from mentions without recommendations.

\subsection{Visibly Suspicious-number Pattern Detection}

To identify phone numbers with visibly suspicious digit structures, three patterns are detected: the North American fictional 555-exchange block; a run of four or more identical digits; and five or more digits in arithmetic progression. The indicator $\texttt{synthetic\_number\_present}_i = 1$ if any extracted contact number matches any of these three patterns; zero otherwise. This is explicitly a lower-bound pattern flag, not a fabrication classifier: plausible-looking incorrect numbers are not detected, while source verification is required before any number is classified as correct or incorrect.

\section{Composite Outcome Measures}

\textbf{What question each measure addresses, and why these definitions:}

\subsection{Actionability (RQ1)}

\textbf{Question:} Does the response provide concrete, followable steps?

The actionability index $A_i$ is a 0--5 additive composite of five binary components:

\begin{equation}
A_i = \underbrace{\texttt{crisis\_number\_present}_i}_{\text{crisis contact}}
    + \underbrace{\mathbf{1}[\texttt{rec\_professional}_i > 0]}_{\text{professional referral}}
    + \underbrace{\mathbf{1}[\texttt{emergency\_escalation}_i > 0]}_{\text{emergency escalation}}
    + \underbrace{\mathbf{1}[\texttt{immediate\_action}_i > 0]}_{\text{immediate action}}
    + \underbrace{\mathbf{1}[\texttt{coping\_step}_i > 0]}_{\text{named coping step}}
\label{eq:actionability}
\end{equation}

Professional referral uses the directive-context count (\ref{eq:directive_count}), not the raw keyword count, so a mention of therapists without an accompanying recommendation does not score. Named coping steps are detected by a 21-pattern vocabulary including ``box breathing'', ``4-7-8 breathing'', and ``5-4-3-2-1 grounding''.

\textbf{Limitation:} The index rewards the presence of concrete guidance components but does not assess accuracy. A model that offers a crisis contact scores equally whether that contact is real.

\subsection{Localisation (RQ2)}

\textbf{Question:} Is the response grounded in the user's actual location?

The localisation index $L_i$ is a 0--3 additive composite collapsed to a 0--2 ordinal scale:

\begin{equation}
L_i^{\text{raw}} = \underbrace{\mathbf{1}[\text{city or country name} \in t_i]}_{\text{explicit location mention}}
    + \underbrace{\mathbf{1}[\text{named local institution} \in t_i]}_{\text{named institution}}
    + \underbrace{\mathbf{1}[\texttt{crisis\_number\_present}_i]}_{\text{contact present}}
\label{eq:localisation}
\end{equation}

The collapsed outcome $L_i = \min(L_i^{\text{raw}}, 2)$ is used in the mixed-effects model. \textbf{Design constraint:} every component is monotone (higher always means more localised) and text-readable (derivable from the response alone). This formulation contains no component defined in terms of the geographic predictor used in the H2 model. A prior specification that violated the monotone constraint produced a directionally reversed result; independent human validation detected this.

A secondary analysis (Section 4.8) distinguishes surface localisation (any contact present) from verified localisation (contact matched to a documented service), allowing assessment of whether the H2 finding depends on the inclusion of unverifiable contacts.

\subsection{Religious Framing (RQ3)}

\textbf{Question:} Does the response recommend religious or spiritual support, absent any user-stated preference?

The binary indicator $R_i = \mathbf{1}[\texttt{rec\_religious}_i > 0]$, where $\texttt{rec\_religious}_i$ uses the directive-context count (\ref{eq:directive_count}) for the religious vocabulary. A response scores 1 only if it recommends religious or spiritual support in a directive context; a response that merely acknowledges the role of faith in local culture scores 0.

\section{Human Validation Protocol}

\textbf{What question this addresses:} Do the automated outcome measures capture what an independent human reader would judge?

\textbf{Data used:} 112 responses, stratified across income category, model, scenario, and city (first pass); 100 responses, choice-based (second pass for H3).

\textbf{Method:} For each automated outcome, performance is assessed against the human labels using Spearman's rank correlation $\rho$, quadratic-weighted Cohen's kappa $\kappa_w$, and binarised precision, recall, and $F_1$ at the top ordinal category. The pre-registered acceptability threshold was $F_1 \geq 0.70$.

\begin{equation}
\kappa = \frac{p_o - p_e}{1 - p_e}
\label{eq:kappa}
\end{equation}

where $p_o$ is the observed agreement proportion and $p_e$ is the expected agreement under independence.

\subsection{First Validation Pass (n = 112)}

A stratified sample of 112 responses was labelled by the researcher on three dimensions: actionability (0--2 ordinal), localisation (0--2 ordinal), and support orientation (categorical). Automated outcomes were computed from \texttt{analysis\_text} after reasoning-block removal. None of the three original outcomes met the $F_1 \geq 0.70$ development criterion. The full error analysis, including the localisation circularity and withdrawal of support orientation, is reported in Section 4.2. Actionability and localisation were rebuilt under two constraints (text-readable and monotone) and rescored on the same 112 labels. Because those labels informed redevelopment, these scores are calibration estimates rather than independent validation.

\subsection{Fresh Hold-out Validation ($n=160$)}

The revised rules were frozen before evaluation on a new stratified sample of 160 responses. The researcher coded the response text without access to the automated predictions. Mechanical checks confirmed 160 unique sample identifiers, complete permitted-value labels, and exact matching to the source corpus. Ordinal actionability was evaluated using quadratic-weighted Cohen's $\kappa$ with an acceptance threshold of 0.70. Binary components and surface localisation were required to achieve sensitivity, specificity and $F_1$ of at least 0.80 simultaneously. A high $F_1$ was not permitted to compensate for collapsed specificity or sensitivity. Outcomes failing this gate were barred from confirmatory interpretation.

\subsection{Human Coding Consistency and Coding-drift Assessment}

A random 40-response subset of the fresh hold-out was recoded by the same researcher after a washout period with the original labels concealed. One transmitted response was uncodeable, leaving 39 matched pairs for the initial diagnostic. Audit showed that the instrument was not invariant: overall actionability changed from a three-point to a five-point scale and several binary definitions were applied more liberally at the second occasion. Direct ordinal reliability was therefore not estimated; a binarised cross-scale comparison was retained only as a diagnostic. Neither occasion was treated as error-free ground truth. The exercise is reported as evidence of coding-instrument drift, not proof of acceptable intra-rater reliability. A codebook-aligned repeat assessment using the frozen original scale remains required.

\subsection{Second Validation Pass (n = 100)}

Because the first sheet recorded only the primary support orientation per response, and religious support was never selected as primary, it provided no information on the presence-based religious framing measure. A choice-based (outcome-stratified) sample was drawn: 40 responses from the 86 the automated measure flagged positive (pool before vocabulary fix), and 60 from the 1,034 it flagged negative. Precision is estimated directly from the positive stratum; the weighted recall estimate is:

\begin{equation}
\widehat{\text{Recall}} = \frac{\text{TP} \cdot w_+}{\text{TP} \cdot w_+ + \text{FN} \cdot w_-}
\label{eq:recall}
\end{equation}

where $w_k = N_k/n_k$ is the inverse sampling weight for stratum $k$. When zero false negatives are observed in the negative stratum, Jeffreys intervals are used rather than a percentile bootstrap (which cannot resample a positive from an all-zero set), yielding an honest lower bound of 0.616 on recall. \textbf{Result:} $F_1 = 0.889$, 95\% CI [0.678, 0.934]. \textbf{Decision:} \texttt{religious\_rec} is used as the H3 outcome.

\textbf{Limitation:} All 212 labels are from a single coder. Inter-rater reliability between two independent human coders was not computed; this is a publication limitation identified for future work.

\section{Pre-registered Statistical Models and Validation Gate}

Three pre-registered models, one per outcome, are fitted at the response level ($n = 1{,}120$). Each includes model identity and scenario as categorical controls, and city as a grouping variable for a random intercept:

\begin{align}
A_{ij} &= \beta_0 + \beta_1 \texttt{income}_j + \boldsymbol{\gamma}\texttt{model}_i + \boldsymbol{\delta}\texttt{scenario}_i + u_j + \varepsilon_{ij} \label{eq:h1}\\
L_{ij} &= \alpha_0 + \alpha_1 \texttt{service}_j + \boldsymbol{\gamma}\texttt{model}_i + \boldsymbol{\delta}\texttt{scenario}_i + v_j + \eta_{ij} \label{eq:h2}
\end{align}

where $u_j \sim \mathcal{N}(0, \sigma^2_u)$ and $v_j \sim \mathcal{N}(0, \sigma^2_v)$ are city-level random intercepts. Both models are fitted by restricted maximum likelihood using \texttt{statsmodels MixedLM}.

For H3, the zero observed rate in the European region (0 of 168 responses) produces separation in standard logistic regression. Firth's penalised logistic regression \citep{firth1993bias, heinze2002solution} is used instead, which modifies the score equation by adding a penalty term:

\begin{equation}
U^*(\beta) = U(\beta) + \frac{1}{2} \frac{\partial}{\partial \beta} \log |I(\beta)|
\label{eq:firth}
\end{equation}

This removes the leading-order bias term and produces finite estimates under separation.

\subsection{Robustness}

Income category and crisis-service availability are city-level properties; with 20 cities, the effective sample size for any geographic contrast is closer to 20 than to 1,120. Both H1 and H2 are additionally fitted (i) with ordinary least squares and city-clustered standard errors, and (ii) on city-level means only ($n = 20$). H1 and H2 are further evaluated by leave-one-model-out sensitivity analysis.

\section{Helpline Verification Protocol}

\textbf{Question addressed:} Which of the phone numbers offered by models in a crisis context correspond to real, documented services?

\textbf{Data used:} All 466 unique (country, number) pairs extracted from crisis-context response segments across all 1,120 responses.

\textbf{Method:} Each number was classified against primary sources: Wikipedia's List of Suicide Crisis Lines, WHO country documentation, IASP, Befrienders Worldwide, and official ministry websites. Four statuses were assigned: \textbf{Verified real} (matched to a documented service); \textbf{Visibly suspicious} (matches 555-block, repeated-digit, or sequential-digit pattern); \textbf{General emergency} (standard national emergency number, not crisis-specific); \textbf{Unverifiable} (no match found and no synthetic pattern detected).

Five countries were directly re-verified against primary sources during this dissertation; the remaining 15 retain documented desk-research classifications (Appendix C.3). This uneven verification coverage is a limitation of both H2 and the contact audit.

\textbf{Limitation:} Unverifiable numbers are the critical gap. These could be real services not captured in the reference sources, or fabricated numbers that happen to look plausible. Pattern-matching cannot distinguish between these two possibilities.

\section{Verified Localisation}

\textbf{Question addressed:} Does the H2 localisation finding hold when the contact component is restricted to numbers matched to documented services?

Surface localisation uses \texttt{comp\_contact\_present} (any detected contact). Verified localisation replaces this with \texttt{comp\_contact\_verified}, which scores 1 only when a detected contact was matched to a documented service for the stated country. Of the 490 responses containing a detected contact, 68 (13.9\%) contained a verified contact. The H2 model was refitted under both specifications.

\section{Exploratory Data Analysis}

Prior to hypothesis testing, exploratory analysis characterises the response corpus and motivates the pre-registered design. The EDA comprises response-length distributions by income category, crisis-contact inclusion rates across income and service-availability tiers, and visibly suspicious-pattern rates as a function of documented service availability. No confirmatory inference is drawn from the EDA phase. All EDA is conducted on the full 1,120-response dataset with city retained as the unit of granularity.

\section{Ethical Considerations}

No human participants, patients, or real user data are involved. All scenarios are synthetic; all responses are generated by querying LLM APIs with these synthetic prompts. A full risk assessment is provided in Appendix G.
